\subsection{哈希表}
线性探测哈希表，使用进程级随机种子的 splitmix64；平均常数复杂度，性能取决于容量和负载。原稿的特定机器测速不能视为通用结论。
\begin{lstlisting}
using u64=unsigned long long;
u64 xorshift(u64 x) {
	static const u64 seed=chrono::steady_clock::now().time_since_epoch().count();
    x += seed+0x9e3779b97f4a7c15;
    x = (x ^ (x >> 30)) * 0xbf58476d1ce4e5b9;
    x = (x ^ (x >> 27)) * 0x94d049bb133111eb;
    return x ^ (x >> 31);
}
template<typename K, typename V, int M=(1<<20)>
struct HashTable{
	static_assert(M>0 && (M&(M-1))==0);
    static const int MASK=M-1;
	int cnt=0;vector<int>e;
	bitset<M> used;
	K k[M]; V v[M];
	void clear(){
		for(auto x:e)used[x]=0;
		e.clear();cnt=0;
	}
	int size(){return cnt;}
	int idx(K x){
		int u=xorshift(x)&MASK;
		for(int step=0;step<M;step++,u=(u+1)&MASK)
            if(!used[u]||k[u]==x)return u;
        return -1; // full table: never probe forever
	}
	V* find(K x){
		int i=idx(x);
		if(i<0||!used[i])return nullptr;
		return &v[i];
	}
	V& operator [](K x){
		int i=idx(x);
        if(i<0)throw overflow_error("HashTable capacity exhausted");
		if(!used[i]){
			used[i]=1;cnt++;
			e.pb(i);
			k[i]=x,v[i]=V{};
		}return v[i];
	}
};
\end{lstlisting}
使用 operator[] 插入或访问，find 返回值指针。此处默认 K、V 为整数；其他键类型需要替换哈希函数，且不提供有序 map 操作。


\paragraph{使用条件与多测} 每组开始调用 \texttt{clear()}；新对象的 bitset 自带零初始化。线性探测仅支持插入、查找和整体清空，不支持直接删槽；建议负载率低于 $0.7$，高负载会退化，最坏单次 $O(M)$。该结构无序，不能代替 map 的有序操作。大对象应放在静态区或堆上。
